% ============================================================
% NIDA Academic Report - English Abstract
% ============================================================

\chapter*{ABSTRACT}
\addcontentsline{toc}{chapter}{ABSTRACT}

\textbf{NIDA: An AI-Powered Child Drawing Analysis System in the Context of House-Tree-Person Test for Psychological Counseling and Guidance}

\vspace{0.5cm}

This study encompasses a digital application project developed to support projective drawing tests used in evaluating children's psychological states in the field of psychological counseling and guidance (PCG) with artificial intelligence technologies. \textbf{NIDA (Neural Imagery Decision Assessment)} is a web-based decision support system that analyzes House-Tree-Person (HTP) drawings of children aged 5-12 and provides academic literature-based interpretations to school psychological counselors and guidance teachers.

The primary aim of this study is to support the process of evaluating children's drawings by psychological counselors in school settings, to provide an assistive tool in detecting early psychological issues, and to contribute to the integration of PCG practices with technology. The system has created a structured Knowledge Base by combining the interpretation rules of the HTP test developed by John Buck (1948) with current academic literature and integrated this knowledge base with Multimodal Large Language Models (MLLM).

NIDA is positioned not as a diagnostic tool but as a decision support system that supports professional evaluation. The system bases each interpretation on academic sources, indicates confidence levels, and transparently presents conflicting academic views. This approach ensures technology integration compatible with the ethical principles and scientific foundations of the PCG field.

This study contributes to the literature on the integration of educational technologies into psychological counseling services, potential uses of artificial intelligence in the PCG field, and the role of digital tools in child assessment.

\vspace{1cm}

\textbf{Keywords:} Psychological counseling and guidance, House-Tree-Person test, projective drawing tests, artificial intelligence, multimodal language models, educational technology, child assessment

\newpage
